% ============================================================
%  SECTION 4: MICROBIOLOGY OF YEASTS IN HONEY
% ============================================================

\section{Microbiology of Yeasts in Honey:
Presence vs.\ Activity vs.\ Adulteration}

The microbiology of honey is a field relatively rarely analyzed
in the context of product authenticity assessment, despite the fact
that microorganisms -- among them osmophilic yeasts -- are an
inseparable component of every honey produced under natural
conditions. This section states and justifies the central
microbiological thesis of the article: the mere presence of
osmophilic yeasts in honey is a biologically normal phenomenon,
does not constitute grounds for declaring the product immature,
let alone adulterated, provided that water activity ($a_w$) and
water content are maintained at levels prescribed by applicable
standards.

%,----------------------------------------------------------
\subsection{Osmophilic yeasts as a natural component of honey}
%,----------------------------------------------------------

\subsubsection{Biological characteristics of osmophilic yeasts}

Osmophilic yeasts, also called osmotolerant or xerotolerant, are
microorganisms capable of growth and metabolism in environments
with very high osmotic pressure, i.e., at water activity $a_w$
close to $0.60$--$0.62$~\cite{ruegg1981,waaw2006}. In honey,
the dominant role is played primarily by species of the genus
\textit{Zygosaccharomyces}:
\textit{Zygosaccharomyces rouxii} (Boutroux) Yarrow, first described
in 1893 as yeasts isolated from sucrose solutions, constitutes the
model species for fungal osmotolerance. Its minimum water activity
for growth is $a_w = 0.62$, making it one of the most resistant
osmotolerant fungi known to
science~\cite{waaw2006}. Under conditions below this threshold,
\textit{Z.~rouxii} remains in a state of cryptobiosis; cells
retain viability and germination capacity but show no metabolic
activity or growth~\cite{waaw2006,ruegg1981}.

More recent studies using culture-dependent and culture-independent
methods (metagenomics) have expanded the list of yeasts isolated
from honey to include further unconventional species, including
\textit{Candida} spp., \textit{Rhodosporidiobolus ruineniae},
\textit{Clavispora lusitaniae}, and \textit{Metschnikowia
chrysoperlae}~\cite{rodriguezmachado2024}.
Metagenomic studies indicate that \textit{Zygosaccharomyces}
remains the dominant genus in both western honey bee
(\textit{Apis mellifera}) honeys and eastern honey bee
(\textit{Apis cerana}) honeys, regardless of country of
origin~\cite{yeastbiodiversity2025}.

\subsubsection{Sources of colonization and prevalence}
Osmophilic yeasts are an integral component of the hive ecosystem
and beekeeping environment, which not only precludes the
justifiability of imposing restrictions in this regard but also
excludes the possibility of their elimination from honey without
applying extreme technological interventions. The main colonization
routes include floral nectar, in which yeast counts are typically
$10^2$--$10^4$~CFU/ml depending on the plant species, microclimate,
and season; pollen and bee bread, which are carriers of active
yeast and lactic acid bacteria populations; as well as the bee's
digestive tract -- the honey crop and intestine constitute a natural
reservoir of microorganisms~\cite{alwaili2012}.
Additional sources include soil, wax, and hive surfaces, whose yeast
biofilm is difficult to eliminate without disinfectants incompatible
with organic production~\cite{alwaili2012}.

The presence of osmophilic yeasts in honey is therefore an inevitable
phenomenon resulting from the biology of the production process,
not from production deficiencies. This is confirmed by literature
from various climatic and geographical zones, encompassing both
western and eastern honey bee
honeys~\cite{yeastbiodiversity2025}.
The volume of yeasts found in normal non-fermenting commercial
honey typically ranges from several dozen to several hundred CFU/g;
values above 1000~CFU/g are recorded in correct samples without
signs of fermentation, provided moisture does not exceed
$18\%$~\cite{analytica2020}.
% #R2.3 #R3: Added specific experimental CFU/g ranges from reviewed studies to provide concrete empirical grounding for the stated typical values
\textcolor{blue}{To provide concrete empirical context: Ruegg and Blanc~\cite{ruegg1981}
reported osmophilic yeast counts of $10$--$10^3$~CFU/g in commercially
traded European honeys showing no fermentation signs; Rodrigue\-z-Machado
et al.~\cite{rodriguezmachado2024} identified counts ranging from below
detection limits to $1.5 \times 10^3$~CFU/g across 120 honey samples
of diverse botanical origin, none of which showed active fermentation
at moisture $\leq 18\%$. Metagenomic surveys confirm that
\textit{Zygosaccharomyces} populations in this CFU/g range represent
a stable, metabolically inactive reservoir under $a_w < 0.60$
conditions~\cite{yeastbiodiversity2025}.}


%,----------------------------------------------------------
\subsection{Water activity as a critical control parameter}
%,----------------------------------------------------------

\subsubsection{Definition and significance of water activity in honey}

Water activity ($a_w$) is a thermodynamic measure of the
availability of water in a food product for chemical reactions
and microbiological activity; defined as the ratio of the vapor
pressure above the product to the vapor pressure of pure water
at saturation at the same
temperature~\cite{waaw2006}.
In contrast to water content (moisture), expressed as mass
percentage, $a_w$ reflects the actual availability of water to
microorganisms and is the primary parameter determining microbial
growth, enzymatic activity, and non-enzymatic browning of the
product~\cite{waaw2006}.

Honey is a supersaturated solution of simple sugars, fructose
($\sim 38\%$) and glucose ($\sim 30\%$), and higher
oligosaccharides, in an environment of minimal water content.
The high osmolarity of the solution resulting from this saturation
causes water molecules to be strongly bound to sugar molecules,
and their thermodynamic activity is drastically
reduced~\cite{zamora2006,waaw2006}. In honey with moisture
$\leq 20\%$, water activity is typically $a_w < 0.60$,
although the exact value depends on sugar composition, water
content, and measurement temperature~\cite{ruegg1981,gleiter2006}.
The study of Ruegg and Blanc~\cite{ruegg1981}, cited by all
subsequent food microbiology textbooks as a fundamental reference
for honey microbiology, demonstrated that for the moisture range
of 16-21\% water activity falls between $a_w = 0.56$ and
$a_w = 0.63$, with the value $a_w = 0.60$ corresponding to
moisture of $\approx 17.5$--$18.5\%$ depending on sugar
composition.

\subsubsection{Osmophilic yeast growth threshold and water activity}

The minimum water activity required for osmophilic yeast growth
is $a_w = 0.61$--$0.62$; below this threshold growth is
completely inhibited and cells remain in a state of anabiosis
without metabolic activity~\cite{ruegg1981,alwaili2012}.
In the context of honey, this means that in a product with moisture
$\leq 18\%$ ($a_w < 0.60$), osmophilic yeasts are biologically
inactive regardless of their count expressed in CFU/g.

The practical implication is direct: 
honey containing even $10^4$~CFU/g of osmophilic yeasts, but with moisture of approximately 17\% and ($a_w \approx 0.57$--$0.58$), should remain microbiologically stable under standard storage conditions~\cite{alwaili2012}.


The Analytica Laboratories document formulates this principle
as a multiparametric model: fermentation is highly probable only
at moisture $> 19\%$ \emph{and} yeast count $> 10$~CFU/g
simultaneously~\cite{analytica2020}, which \textit{a~contrario}
precludes product disqualification on the basis of the CFU/g
parameter alone.

For reasons of microbiological stability, water activity is an
incomparably more important parameter than the percentage water
content. It should be noted that for certain types of honey
(e.g., heather honey) legal requirements permit water content
levels above 20\%. This is due to the fact that the specific
composition of these honeys ensures microbiological stability at
water content exceeding 20\%, and these honeys should simultaneously
be considered ``mature'' at the values specified in legal
requirements. If water activity studies demonstrate that the water
activity of these honeys does not exceed the threshold
($a_w < 0.60$), this parameter would constitute a significantly
more universal parameter for determining the microbiological
stability and degree of maturity of honey.

\subsubsection{Glucose crystallization as a dynamic factor modifying water activity}

The natural process of glucose crystallization in honey is a factor
deserving particular attention in the context of fermentation risk,
as it constitutes a mechanism by which initially microbiologically
stable honey can become susceptible to fermentation without any
external intervention. During crystallization, glucose transitions
from the liquid phase to the solid phase as glucose monohydrate
(C$_6$H$_{12}$O$_6 \cdot$ H$_2$O). Each glucose molecule binds
one water molecule, effectively removing it from the liquid
phase~\cite{zamora2006}. However, the liquid phase remaining
after glucose crystallization becomes enriched in fructose and
water, causing an \emph{increase} in water activity in the
remaining liquid phase, even though the total water content in
the sample does not change~\cite{zamora2006,ijfans2019}.

Zamora and Chirife~\cite{zamora2006} in a study of 50 samples
of crystallized and re-dissolved Argentine honeys demonstrated
that crystallization causes an increase in the $a_w$ of the liquid
phase by $0.03$--$0.04$ units compared to honey in the liquid
state. For honeys with initial moisture of $17.5$--$18.0\%$
($a_w \approx 0.58$--$0.59$), crystallization can shift the
$a_w$ of the liquid phase above the threshold of $0.62$, thereby
activating the osmophilic yeast population present in the
product~\cite{zamora2006,ijfans2019}.
This phenomenon leads to a key methodological conclusion: the
microbiological control of honey must take into account the phase
state of the product (liquid vs.\ crystallized), since the same
CFU/g count carries different fermentation risk depending on whether
the honey is liquid or crystallized. These processes further support the need to interpret yeast counts together with water activity and the physical state of honey.

Given that honey crystallization proceeds in a differentiated
manner (i.e., through different crystal forms with regard to both
structure and dynamics of formation), changes in water activity
and thus the potential for yeast proliferation in crystallizing
honey are variable. This variability is observed not only with
respect to the time/progress of crystallization but also the
sampling point within the container. Considering that crystallization
propagation occurs at the molecular level, the greatest potential
for water concentration occurs in the pericrystalline and
intercrystalline spaces. Mainly through diffusion, this fluctuation
tends to equalize in areas increasingly distant from the crystal
surface. It should be noted that the interphase equilibrium is a
dynamic process in which the direction of phase transitions is
determined by temperature (variable storage conditions). Depending
on the conditions accompanying these phase transitions (such as
temperature, solution density, oxygen availability, honey aeration,
the free spaces formed in honey during crystallization, the rate of
crystallization resulting from the proportions of individual sugars,
etc.), the potential for yeast proliferation in solution is dynamic
and varies throughout the volume of honey. A CFU/g measurement
without precise specification of the phase state and~$a_w$ is
therefore methodologically incomplete.

From the above it follows that if water activity could determine
honey maturity, if this parameter were specified in legal
requirements, the water activity value should be set at a level
achievable at the time of honey extraction, and should also be
maintained (without further technical measures) despite differences
in temperature and phase transitions occurring during storage and
distribution.

%,----------------------------------------------------------
\subsection{Yeast count as a risk indicator, not a disqualifying parameter}
%,----------------------------------------------------------

\subsubsection{Multiparametric model of fermentation risk}

Analysis of the scientific literature and practical documents
indicates the widespread acceptance of a multiparametric model
of honey fermentation risk, in which no single indicator is
sufficient to determine the microbiological stability or
instability of the product. Fermentation risk is a function
of at least four variables:

\begin{equation}
R_f = f\!\left(a_w,\; T,\; t,\; N_{CFU}\right)
\label{eq:fermentation_risk}
\end{equation}

\noindent where: $R_f$, fermentation risk (probability of initiating
active fermentation), $a_w$, water activity, $T$, storage
temperature [$^\circ$C], $t$, storage time [days],
$N_{CFU}$, osmophilic yeast count [CFU/g].

According to this model, yeasts can ferment honey only when
\emph{all} growth-promoting parameters are simultaneously fulfilled:
$a_w$ must exceed the threshold of $0.62$, temperature must fall
within the optimal range for yeasts ($10$--$27\,^\circ$C),
and the exposure time must be sufficient to reach the cell density
initiating noticeable fermentation~\cite{apinz2021}.
The New Zealand Beekeeping Industry document states directly that
honey with moisture $> 19\%$ and yeast count $> 10$~CFU/g is
susceptible to fermentation, which \emph{a contrario} implies
that honey with moisture $< 17\%$, even with a yeast count many
times higher, does not ferment~\cite{apinz2021}.

Table~\ref{tab:fermentation_risk} illustrates the matrix model
of fermentation risk as a function of moisture and yeast count,
based on available empirical data.

\begin{table}[ht]
\begin{adjustwidth}{-\extralength}{0cm}
\centering
\caption{Estimated honey fermentation risk as a function
         of moisture and osmophilic yeast count
         (storage temperature $10$--$27\,^\circ$C).
         Source: own compilation based on~\cite{analytica2020,apinz2021}.}
\label{tab:fermentation_risk}
\small
\begin{tabularx}{\fulllength}{l *{3}{>{\centering\arraybackslash}X}}
\hline
\multirow{2}{*}{\textbf{Moisture [\%]}} &
\multicolumn{3}{c}{\textbf{Yeast count [CFU/g]}} \\
\cline{2-4}
 & $< 10$ & $10$--$1000$ & $> 1000$ \\
\hline
$\leq 17.0$       & no risk       & no risk               & no risk \\
$17.1$--$18.0$  & no risk       & low risk              & moderate risk \\
$18.1$--$19.0$  & no risk       & moderate risk         & high risk \\
$19.1$--$20.0$  & low risk      & high risk             & fermentation very probable \\
$> 20.0$          & moderate risk & fermentation probable & fermentation inevitable \\
\hline
\end{tabularx}
\end{adjustwidth}
\end{table}



% #R1.4: Added note clarifying heuristic nature of Table 2 risk categories — "no risk" and "inevitable fermentation" are estimates from two practitioner documents, not experimentally validated thresholds
\textcolor{blue}{The risk categories presented in Table~\ref{tab:fermentation_risk} are heuristic estimates compiled from two practitioner-oriented documents~\cite{analytica2020,apinz2021} and represent the consensus of applied honey quality practice. They have not been derived from a controlled experimental study under defined storage conditions; the term ``no risk'' should be understood as indicating negligible probability of fermentation initiation under normal commercial storage, not an absolute thermodynamic guarantee. The table is intended as a decision-support tool and should be interpreted alongside direct $a_w$ measurement rather than used as a standalone classification system.}

\subsubsection{Limitations of the plate method (CFU/g) as a risk assessment tool}

% The standard method for determining yeast and mold counts in food
% is based on plating dilutions of the product onto selective medium
% (e.g., Dichloran Rose Bengal Chloramphenicol medium, DRBC Agar,
% ISO 21527-1:2008) and colony counting after incubation at
% $25\,^\circ$C for 5~days. 

The standard enumeration of yeasts and molds in foods is commonly based on plating serial dilutions on selective media, followed by incubation and colony counting.


This method, while widely used and well-validated, has a number of significant
limitations in the context of honey fermentation risk assessment:


\begin{enumerate}
    \item \textbf{Failure to distinguish viable from dead cells:}
          the plate method counts only cells capable of forming
          colonies on solid medium under incubation conditions;
          living cells in a state of deep anabiosis induced by
          low honey $a_w$ may fail to grow on the medium or grow
          atypically, leading to underestimation of the true yeast
          count;
    \item \textbf{Lack of correlation with in situ $a_w$:} the
          CFU/g result is a measure of the potential population size,
          but does not indicate whether under the conditions prevailing
          in the tested product this population is active;
          only the combination of CFU/g with $a_w$ measurement allows
          assessment of the actual fermentation
          risk~\cite{ruegg1981};
    \item \textbf{Lack of species specificity:} standard selective
          media do not differentiate yeast species with different
          fermentation potential; \textit{Metschnikowia}
          spp.\ with low fermentation potential are counted
          identically to \textit{Z.~rouxii} with high
          potential~\cite{rodriguezmachado2024}.
\end{enumerate}

\textbf{Legal-scientific conclusion:} in light of the above
methodological limitations, the mere determination of yeast count
(CFU/g) in honey, without a simultaneous measurement of $a_w$,
determination of the product's phase state, storage temperature,
and time since harvest, does not constitute a sufficient scientific
basis for assessing fermentation risk, let alone for disqualifying
the product from commercial trade. Given that no applicable standard
(Codex, EU, ISO) establishes a CFU/g limit for honey
(cf.~Section~2), a potential inspector's decision to disqualify
a product based solely on CFU/g count would be devoid of both
scientific and legal grounds.

%,----------------------------------------------------------
\subsection{Fermentation as a dynamic and post-harvest event}
%,----------------------------------------------------------

The Apimondia position and much of the literature on honey
adulteration tacitly assumes that fermentation detected in honey
is evidence of its immaturity at the time of harvest.
This assumption is, however, incorrect from the perspective of
microbiological dynamics, since fermentation is an event that can
occur at any stage of the supply chain, not only at the time
the product leaves the hive~\cite{apinz2021,analytica2020}.

Correctly harvested honey that meets all normative criteria,
with moisture of, for example, $18.5\%$ ($a_w \approx 0.60$)
and a yeast count of $200$~CFU/g, stored for six months in a
warehouse at a temperature of $22\,^\circ$C and relative humidity
above 70\%, may locally exhibit higher water activity
(crystallization, temperature changes, air access, etc.), shifting
the $a_w$ above the yeast growth threshold and initiating
fermentation~\cite{apinz2021}.

Border inspection, sampling imported honey after several months
of transport and storage, takes measurements at a point removed
in time and space from the moment of harvest; the determination of
fermentation or elevated active yeast count at that moment does not
allow, without analysis of temperature and humidity documentation
in the supply chain (\textit{cold chain monitoring}), reliable
reconstruction of the product's state at the time of honey
extraction (verifying the degree of capping) or departure from
the apiary~\cite{tsagkaris2021,schoder2026}.
Therefore, fermentation observed at the point of inspection should not be treated, by itself, as proof of deliberate adulteration at the production stage. It may indicate product instability or non-compliance, but attribution to intentional immature honey production requires additional evidence from the extraction, processing, storage, or supply-chain stages.

Deliberate adulteration consisting of
harvesting immature honey followed by technological treatments
aimed at its microbiological stabilization can only be established
through controls at the honey extraction stage, and effective
limitation may be achieved by specifying prohibited technologies
(e.g., vacuum dehydration). The non-use of prohibited technologies
should be a mandatory stage of official controls of apiaries and
establishments involved in honey extraction and placing on the
market. Neither yeast count nor fermentation permits inference
about this.
